\chapter{Estado del Arte y Fundamentos Teóricos}
\label{ch:estado-arte}

\section{Grandes modelos de lenguaje}

\subsection{Arquitectura Transformer}

Hasta 2017, los modelos dominantes en procesamiento del lenguaje natural eran las redes
recurrentes, que procesan el texto token a token acumulando el contexto en un vector de estado
oculto. No obstante, esa cadena de multiplicaciones hace que los gradientes se desvanezcan o
exploten durante el entrenamiento, dificultando el aprendizaje de dependencias entre palabras
lejanas\footnote{%
  El problema del gradiente desvaneciente fue estudiado por Hochreiter~(1991) y formalizado por
  Bengio et al.~(1994). LSTM y GRU paliaron el síntoma introduciendo compuertas, pero no
  eliminaron la causalidad secuencial.%
}. La propuesta de Vaswani et al.~\cite{vaswani2017attention} fue eliminar la recurrencia por
completo: en lugar de propagar el contexto de forma secuencial, cada token puede relacionarse
directamente con cualquier otro de la secuencia en un único paso paralelizable.

El mecanismo que lo hace posible es la \textbf{atención}. A partir de la entrada se derivan,
mediante proyecciones lineales de la matriz de embeddings, tres matrices: $Q$ (\textit{queries}),
que representa qué información busca cada posición; $K$ (\textit{keys}), que representa qué
información ofrece cada posición para ser encontrada; y $V$ (\textit{values}), que contiene el
contenido que se recupera cuando una posición es seleccionada. La atención se calcula como:

\begin{equation}
  \text{Attention}(Q, K, V) = \text{softmax}\!\left(\frac{QK^\top}{\sqrt{d_k}}\right)V
  \label{eq:attention}
\end{equation}

El producto $QK^\top$ mide la afinidad entre cada par de posiciones; dividirlo por $\sqrt{d_k}$
evita que los valores crezcan tanto que el \textit{softmax} sature y los gradientes desaparezcan.
El resultado es una suma ponderada de los valores $V$, donde los pesos reflejan cuánto «presta
atención» cada posición a las demás.

Una sola aplicación de esta operación captura un tipo de relación a la vez. Por eso la
arquitectura usa \textit{multi-head attention}: aplica la ecuación~\eqref{eq:attention}
en paralelo en múltiples subespacios proyectados independientemente, cada uno libre de especializarse
en un patrón distinto —sintáctico, semántico, correferencial— y concatena las salidas al final.
Así, el modelo puede razonar simultáneamente sobre estructura local y dependencias globales.

La arquitectura original combina dos bloques con roles complementarios. El \textbf{codificador}
(\textit{encoder}) procesa la secuencia de entrada completa con atención bidireccional,
produciendo una representación contextualizada de cada token; es la elección natural para tareas
de comprensión. El \textbf{decodificador} (\textit{decoder}) genera la salida token a token de
forma autoregresiva, usando atención enmascarada para no ver posiciones futuras y atención cruzada
sobre la salida del codificador para anclar la generación en la entrada.

Ambos bloques inyectan codificaciones posicionales sinusoidales para preservar el orden sin recurrir
a la recurrencia. Escalada en datos y parámetros, esta arquitectura dio lugar a los grandes modelos
de lenguaje actuales.

\subsection{Modelos de generación de texto}

Los grandes modelos de lenguaje para generación de texto adoptan una variante \textit{decoder-only}
del Transformer: eliminan el bloque codificador y emplean únicamente el decodificador con atención
causal sobre toda la secuencia. Esta simplificación concentra el entrenamiento en un único objetivo
—predecir el siguiente token sobre el corpus completo— y ha demostrado escalar de forma
especialmente eficiente a cientos de miles de millones de parámetros. Modelos como GPT-5, Claude
o Gemini responden a esta arquitectura.

Un hallazgo determinante sobre estos modelos llegó con Brown et al.~\cite{brown2020gpt3}: los modelos
con suficiente capacidad exhiben \textit{few-shot learning}, es decir, dados unos pocos ejemplos en
el propio prompt son capaces de generalizar la tarea sin actualizar sus pesos. Esto redujo drásticamente
la dependencia de conjuntos de datos etiquetados para cada aplicación específica y abrió la puerta a
sistemas adaptables sin \textit{fine-tuning}.

El siguiente paso fue el \textit{fine-tuning}. Los modelos preentrenados son competentes, pero no
necesariamente útiles: pueden generar texto coherente sin seguir instrucciones ni respetar
preferencias humanas. El aprendizaje por refuerzo a partir de retroalimentación humana
(\textit{RLHF}) resolvió esto en dos fases: primero, evaluadores humanos clasifican pares de
respuestas del modelo, entrenando un modelo de recompensa que aprende a predecir qué respuesta
preferiría un humano; después, el modelo de lenguaje se ajusta mediante aprendizaje por refuerzo
para maximizar esa recompensa. El resultado es un modelo que sigue instrucciones, reconoce cuando
no sabe algo y produce salidas estructuradas con mucha mayor fiabilidad. Son estos modelos —
preentrenados a escala y alineados con instrucciones— los que \textbf{Seshat} emplea como base
de sus agentes.

\subsection{\textit{Fine-tuning} y adaptación mediante prompts}

El \textit{fine-tuning} permite especializar un modelo preentrenado en un dominio concreto
mediante entrenamiento adicional sobre datos etiquetados. Aunque eficaz, introduce costes de
datos, cómputo y mantenimiento que lo hacen poco adecuado para sistemas que deben adaptarse a
cambios frecuentes de modelo base o de dominio.

Una alternativa más ligera es el \textit{prompt engineering}, que explota la capacidad de los
\ac{LLM} de generalizar a partir del contexto sin modificar sus pesos. Más allá de proporcionar
ejemplos de salida esperada (\textit{few-shot learning}), el diseño cuidadoso de prompts para
tareas de extracción estructurada combina varias técnicas complementarias: el razonamiento en
cadena~\cite{wei2022cot} descompone la tarea en pasos intermedios verificables antes de emitir
la respuesta; el anclaje en cita literal obliga al modelo a localizar el fragmento de texto que
respalda cada nodo antes de describirlo, reduciendo la alucinación; las compuertas lógicas
codifican condiciones de inclusión y exclusión que el modelo debe evaluar explícitamente; y los
contraejemplos etiquetados delimitan fronteras semánticas entre categorías próximas. La salida
estructurada forzada mediante esquemas Pydantic elimina la ambigüedad de formato y habilita
reintentos automáticos ante violaciones de tipo. \textbf{Seshat} opta por este segundo enfoque;
las técnicas concretas se describen en el anexo~\ref{anx:tecnicas-prompts}.

\section{Recuperación aumentada por generación}

Los \ac{LLM} tienen un límite estructural: su conocimiento queda congelado en el momento del
entrenamiento. La generación aumentada por recuperación (\ac{RAG})~\cite{lewis2020rag} aborda
este problema combinando el modelo generativo con un módulo de recuperación: antes de generar,
el sistema consulta una base de conocimiento externa, selecciona los fragmentos más relevantes e
incluye estos en el prompt. La respuesta queda así anclada en evidencia concreta y el
conocimiento del sistema puede actualizarse sin re-entrenar el modelo.

La pieza central del módulo de recuperación son los \textit{embeddings}: representaciones
vectoriales densas que codifican el significado semántico de un texto en un espacio de alta
dimensión, de modo que textos con significado similar queden cerca en ese espacio. La semejanza
entre vectores, medida típicamente mediante similitud coseno, permite identificar los fragmentos
más relevantes para una consulta dada. Bases de datos vectoriales como \texttt{pgvector}
proporcionan la infraestructura de indexación y búsqueda eficiente sobre estas representaciones.

El paradigma ha evolucionado más allá de la búsqueda plana. Para mejorar la cobertura, se han
propuesto estrategias de expansión de consulta: en lugar de buscar con la pregunta original,
el sistema genera múltiples variantes con un LLM y fusiona los resultados, diversificando los
candidatos recuperados~\cite{gao2024ragsurvey}. En la etapa de selección, el reranking con
modelos \textit{cross-encoder}~\cite{nogueira2020reranking} mejora la precisión al comparar
directamente la consulta con cada candidato, en lugar de aproximar la relevancia con la similitud
coseno. GraphRAG~\cite{edge2025graphrag} va más allá y construye un grafo de entidades a partir
de un corpus documental para enriquecer la recuperación con contexto estructural. \textbf{Seshat}
incorpora las dos primeras estrategias —multi-query y reranking opcional— e invierte la premisa
de GraphRAG: usa \ac{RAG} no para responder preguntas, sino para construir y mantener el propio
grafo, recuperando los nodos preexistentes más relevantes para determinar qué relaciones semánticas
introduce cada nuevo nodo extraído.

\section{Procesamiento del habla}

\subsection{Reconocimiento automático del habla}

Hasta mediados de la década pasada, los sistemas de reconocimiento automático del habla (\ac{ASR})
se construían como cadenas de módulos especializados: un modelo acústico que asignaba
probabilidades a fonemas, un modelo de lenguaje que los combinaba en palabras y un decodificador
que buscaba la secuencia más probable. El aprendizaje profundo fue reemplazando estos módulos por
redes neuronales entrenadas conjuntamente, y la arquitectura Transformer completó la transición
hacia modelos encoder-decoder que aprenden directamente la correspondencia entre audio y texto sin
alineamiento explícito.

Whisper~\cite{radford2023whisper} demostró hasta dónde podía llegar este enfoque: entrenado sobre
680.000 horas de audio con transcripciones débilmente supervisadas obtenidas de internet, produce
un modelo robusto y multilingüe que no requiere \textit{fine-tuning} para nuevos dominios. Su
codificador procesa espectrogramas de log-Mel y su decodificador genera la transcripción de forma
autoregresiva.

\textbf{Seshat} delega la transcripción en AssemblyAI, un servicio de la misma familia. La
interfaz de transcripción es abstracta y sustituible sin modificar la lógica del \textit{pipeline}.

\subsection{Síntesis de voz}

La síntesis de voz (\textit{Text-to-Speech}, TTS) siguió un arco similar. Los sistemas clásicos
producían voz inteligible pero perceptiblemente artificial. WaveNet~\cite{oord2016wavenet} fue el
primer modelo neuronal capaz de generar formas de onda con naturalidad comparable a la voz humana,
aunque con un coste computacional prohibitivo para tiempo real. Los modelos posteriores separaron
la generación del espectrograma de mel\footnote{%
  Un espectrograma de mel representa cómo varía en el tiempo la energía de una señal de audio en
  distintas bandas de frecuencia, con una escala (\textit{mel}) que espacia esas bandas de forma
  similar a como el oído humano percibe el tono, en lugar de linealmente en Hz.%
} de la síntesis de la forma de onda, reduciendo la latencia.
VITS~\cite{kim2021vits} unificó ambas etapas: combina un  autocodificador variacional con un modelo
de flujo normalizado y entrenamiento adversarial para generar audio de alta calidad en una sola pasada.

En este trabajo, la TTS no forma parte del \textit{pipeline} en tiempo de ejecución. ElevenLabs,
un servicio comercial de alta fidelidad, se empleó exclusivamente fuera de línea para generar
grabaciones sintéticas de reuniones multi-locutor —con ruido de fondo de oficina opcional— que
ejercitan la ruta de transcripción e ingesta de audio de extremo a extremo.

\section{Consideraciones éticas y legales}

Las transcripciones de reuniones técnicas pueden contener información sensible: datos personales,
decisiones estratégicas confidenciales o referencias a terceros. El Reglamento General de
Protección de Datos (RGPD) exige que el tratamiento de datos personales cuente con una base
legal, que los interesados sean informados y que los datos se almacenen de forma segura.
\textbf{Seshat} no implementa en su versión actual ningún mecanismo de detección ni anonimización
de información personal identificable (PII); esta capacidad queda identificada como línea de
trabajo futuro y se aborda en el capítulo~\ref{ch:discusion}.

Un aspecto igualmente relevante es el consentimiento informado de los participantes: más allá del
cumplimiento del RGPD, la existencia de un sistema que estructura permanentemente el contenido de
las reuniones puede desincentivar la comunicación abierta si los participantes no son conscientes
de su funcionamiento. Para uso interno, esto es una cuestión de confianza organizativa tanto como
de obligación legal.

Más allá de la privacidad y el consentimiento, los \ac{LLM} introducen un riesgo propio: las
alucinaciones, es decir, la generación de afirmaciones falsas pero plausibles. Estas pueden dar
lugar a nodos incorrectos en la base de conocimiento, con consecuencias reales si esos nodos
informan decisiones futuras. \textbf{Seshat} mitiga este riesgo mediante la revisión humana de
los nodos de baja confianza, el anclaje de cada nodo a una cita literal de la transcripción y el
registro completo de prompts e inferencias en MLflow~\cite{zaharia2018mlflow} para auditoría.

Un tercer vector de riesgo es la inyección de prompt: un actor que controle el contenido de la
transcripción podría insertar instrucciones que manipulen el comportamiento del agente.
\textbf{Seshat} lo previene mediante delimitadores explícitos que separan el contenido procesado
de las instrucciones del sistema, impidiendo que el texto de entrada pueda reinterpretar el rol
del agente.
